\documentclass[a4paper,10pt]{article}

\usepackage[utf8]{inputenc}
\usepackage[T1]{fontenc}
\usepackage{geometry}
\usepackage{booktabs}
\usepackage{array}
\usepackage{longtable}
\usepackage{listings}
\usepackage[hidelinks]{hyperref}

\geometry{
    top=1.8cm,
    bottom=1.8cm,
    left=1.8cm,
    right=1.8cm
}

\setlength{\parindent}{0pt}
\setlength{\parskip}{4pt}

\lstset{
    basicstyle=\ttfamily\small,
    frame=single,
    breaklines=true,
    columns=fullflexible
}

\title{\textbf{DEOS BMS Service Protocol}\\
\large CAN-FD ICD Tamamlayıcı Dokümanı}
\author{DEOS -- Dokuz Eylül Otonom Sistemler}
\date{}

\begin{document}

\maketitle

\section{Amaç ve Kapsam}

Bu doküman, ana \textit{DEOS CAN-FD Interface Control Document}
içerisinde tanımlanan BMS haberleşme yapısını tamamlamak amacıyla
hazırlanmıştır.

Doküman; Textual Node, Ground Control Node, Main STM32,
micro-ROS Node veya diğer yetkili DEOS node'larının batarya yönetim
sisteminden bilgi alırken kullanacağı servis, komut ve payload
yapılarını tanımlar.

Ana ICD içerisinde tanımlanan aşağıdaki yapılar bu dokümanda tekrar
tanımlanmaz:

\begin{itemize}[nosep]
    \item 29-bit CAN-FD Extended Identifier yapısı,
    \item Priority alanı,
    \item Source ve Destination Node alanları,
    \item Message Class yapısı,
    \item genel routing mekanizması,
    \item Protocol Version ve Sequence alanları,
    \item CAN-FD DLC ve padding kuralları,
    \item Little-Endian byte-order kuralı.
\end{itemize}

Bu doküman ana ICD ile birlikte kullanılmalıdır.

% =========================================================
\section{BMS Sistemindeki Node'lar}

BMS sistemi aşağıdaki DEOS node'larından oluşur:

\begin{center}
\begin{tabular}{c l l}
\toprule
\textbf{Node ID} & \textbf{Symbol} & \textbf{Görev} \\
\midrule

0x20 &
DEOS\_NODE\_BMS\_MAIN &
Ana batarya BMS \\

0x21 &
DEOS\_NODE\_BMS\_AUX &
Yardımcı batarya BMS \\

0x30 &
DEOS\_NODE\_BMS\_GATEWAY &
Classic CAN -- CAN-FD gateway \\

\bottomrule
\end{tabular}
\end{center}

BMS\_MAIN ve BMS\_AUX fiziksel olarak BMS Gateway arkasında
bulunsa dahi DEOS protokolünde bağımsız mantıksal node'lar olarak
adreslenir.

Gateway yalnızca SYSTEM ve DIAGNOSTIC servislerini kendi adına
destekler.

\textbf{BMS Service, BMS Gateway tarafından yerel servis olarak
desteklenmez.}

BMS Service mesajları doğrudan:

\begin{itemize}[nosep]
    \item DEOS\_NODE\_BMS\_MAIN veya
    \item DEOS\_NODE\_BMS\_AUX
\end{itemize}

adreslerine gönderilmelidir.

Gateway gerekli Classic CAN--CAN-FD protokol dönüşümünü gerçekleştirir
ancak üst seviye DEOS mesajında mantıksal BMS kimliği korunur.

Örneğin BMS\_MAIN verisi gateway üzerinden taşınsa bile:

\begin{lstlisting}
Source = DEOS_NODE_BMS_MAIN
\end{lstlisting}

olarak gönderilir.

Gateway'in kendi mesajlarında ise:

\begin{lstlisting}
Source = DEOS_NODE_BMS_GATEWAY
\end{lstlisting}

kullanılır.

% =========================================================
\section{Node--Service İlişkisi}

BMS alt sistemi içerisinde servis sahipliği aşağıdaki gibidir:

\begin{center}
\begin{tabular}{l c c c}
\toprule
\textbf{Node} &
\textbf{SYSTEM} &
\textbf{BMS} &
\textbf{DIAGNOSTIC} \\
\midrule

BMS\_MAIN    & X & X & X \\
BMS\_AUX     & X & X & X \\
BMS\_GATEWAY & X & -- & X \\

\bottomrule
\end{tabular}
\end{center}

Bu tablo ilgili servisin node üzerinde yerel olarak işlendiğini
gösterir.

% =========================================================
\section{BMS Service}

\subsection{Genel}

BMS Service ID:

\begin{lstlisting}
DEOS_SERVICE_BMS = 0x04
\end{lstlisting}

olarak tanımlıdır.

BMS Service yalnızca BMS\_MAIN ve BMS\_AUX node'ları için geçerlidir.

Serviste iki farklı veri erişim modeli kullanılır:

\begin{enumerate}[nosep]
    \item Genel BMS bilgilerinin tek bir STATUS mesajında alınması,
    \item Değişken sayıda bulunan verilerin liste şeklinde alınması.
\end{enumerate}

Hücre voltajları ve sıcaklık sensörleri gibi eleman sayısı batarya
konfigürasyonuna göre değişebilen veriler liste tipi mesajlarla
aktarılır.

\subsection{BMS Command Registry}

\begin{center}
\begin{tabular}{c l}
\toprule
\textbf{ID} & \textbf{Command} \\
\midrule

0x01 & DEOS\_CMD\_BMS\_GET\_STATUS \\
0x02 & DEOS\_CMD\_BMS\_STATUS \\

0x03 & DEOS\_CMD\_BMS\_GET\_CELL\_VOLTAGES \\
0x04 & DEOS\_CMD\_BMS\_CELL\_VOLTAGE \\

0x05 & DEOS\_CMD\_BMS\_GET\_TEMPERATURES \\
0x06 & DEOS\_CMD\_BMS\_TEMPERATURE \\

0x07 & DEOS\_CMD\_BMS\_GET\_BALANCING \\
0x08 & DEOS\_CMD\_BMS\_BALANCING \\

\bottomrule
\end{tabular}
\end{center}

% =========================================================
\section{Genel BMS Durum Bilgisi}

\subsection{GET\_STATUS}

Bir uygulama BMS'in temel çalışma bilgilerini almak istediğinde:

\begin{lstlisting}
Message Class = DEOS_CLASS_REQUEST
Service       = DEOS_SERVICE_BMS
Command       = DEOS_CMD_BMS_GET_STATUS
\end{lstlisting}

mesajını gönderir.

Command-specific request payload bulunmaz.

Örnek:

\begin{lstlisting}
Source        = DEOS_NODE_GROUND_CONTROL
Destination   = DEOS_NODE_BMS_MAIN
Message Class = DEOS_CLASS_REQUEST
Service       = DEOS_SERVICE_BMS
Command       = DEOS_CMD_BMS_GET_STATUS
\end{lstlisting}

\subsection{STATUS}

GET\_STATUS isteğine:

\begin{lstlisting}
Message Class = DEOS_CLASS_RESPONSE
Service       = DEOS_SERVICE_BMS
Command       = DEOS_CMD_BMS_STATUS
\end{lstlisting}

mesajı ile cevap verilir.

STATUS mesajı BMS'in genel durumunu tek frame içerisinde özetler.

Command payload aşağıdaki şekilde tanımlanır:

\begin{center}
\begin{tabular}{c c c l}
\toprule
\textbf{Offset} &
\textbf{Boyut} &
\textbf{Tip} &
\textbf{Alan} \\
\midrule

0  & 2 byte & uint16\_t & Pack Voltage \\
2  & 2 byte & int16\_t  & Pack Current \\
4  & 2 byte & uint16\_t & State of Charge \\
6  & 2 byte & int16\_t  & BMS Temperature \\
8  & 2 byte & uint16\_t & Minimum Cell Voltage \\
10 & 2 byte & uint16\_t & Maximum Cell Voltage \\
12 & 2 byte & uint16\_t & Nominal Pack Voltage \\
14 & 1 byte & uint8\_t  & Cell Count \\
15 & 1 byte & uint8\_t  & Temperature Sensor Count \\

\bottomrule
\end{tabular}
\end{center}

Ölçekler:

\begin{center}
\begin{tabular}{l c c}
\toprule
\textbf{Alan} & \textbf{Scale} & \textbf{Birim} \\
\midrule

Pack Voltage          & 0.01 & V \\
Pack Current          & 0.1  & A \\
State of Charge       & 0.01 & \% \\
BMS Temperature       & 0.1  & $^\circ$C \\
Minimum Cell Voltage  & 0.001 & V \\
Maximum Cell Voltage  & 0.001 & V \\
Nominal Pack Voltage  & 0.01 & V \\

\bottomrule
\end{tabular}
\end{center}

Nominal Pack Voltage bilgisi kullanılan BMS tarafından sağlanmıyorsa
ilgili alan \texttt{0x0000} olarak gönderilebilir.

Cell Count ve Temperature Sensor Count alanları bilgilendirme amacıyla
bulunur. Liste aktarımı bu değerlere bağımlı değildir; liste sonunun
belirlenmesinde End Of List mesajı kullanılır.

% =========================================================
\section{Liste Tipi Veri Aktarımı}

Hücre voltajları, sıcaklık sensörleri ve benzeri değişken uzunluklu
veriler için ortak liste aktarım yöntemi kullanılır.

Bir liste request'i alındığında hedef BMS:

\begin{enumerate}[nosep]
    \item Mevcut tüm elemanları sırayla gönderir.
    \item Her CAN-FD frame'i yalnızca bir eleman içerir.
    \item Ardışık elemanlar nominal 100 ms aralıkla gönderilir.
    \item Son elemandan sonra End Of List mesajı gönderilir.
\end{enumerate}

İlk elemanın gönderilmesi için 100 ms beklenmesi zorunlu değildir.

Liste transferlerinde index değeri:

\begin{center}
\begin{tabular}{c l}
\toprule
\textbf{Index} & \textbf{Anlam} \\
\midrule
0x00 & END\_OF\_LIST \\
0x01--0xFF & Gerçek eleman numarası \\
\bottomrule
\end{tabular}
\end{center}

Dolayısıyla gerçek hücre ve sensör numaralandırması 1'den başlar.

Liste uzunluğunun request öncesinde bilinmesi gerekmez.

% =========================================================
\section{Hücre Voltajlarının Okunması}

\subsection{GET\_CELL\_VOLTAGES}

Tüm hücre voltajlarını istemek için:

\begin{lstlisting}
Message Class = DEOS_CLASS_REQUEST
Service       = DEOS_SERVICE_BMS
Command       = DEOS_CMD_BMS_GET_CELL_VOLTAGES
\end{lstlisting}

kullanılır.

Request command-specific payload içermez.

\subsection{CELL\_VOLTAGE}

Her hücre aşağıdaki mesaj ile ayrı ayrı gönderilir:

\begin{lstlisting}
Message Class = DEOS_CLASS_RESPONSE
Service       = DEOS_SERVICE_BMS
Command       = DEOS_CMD_BMS_CELL_VOLTAGE
\end{lstlisting}

Command payload:

\begin{center}
\begin{tabular}{c c c l}
\toprule
\textbf{Offset} &
\textbf{Boyut} &
\textbf{Tip} &
\textbf{Alan} \\
\midrule

0 & 1 byte & uint8\_t  & Cell Index \\
1 & 2 byte & uint16\_t & Cell Voltage \\

\bottomrule
\end{tabular}
\end{center}

Cell Voltage ölçeği:

\begin{lstlisting}
Physical Voltage = Raw Value * 0.001 V
\end{lstlisting}

Örnek:

\begin{lstlisting}
Cell Index   = 0x03
Cell Voltage = 3652

-> Cell 3 = 3.652 V
\end{lstlisting}

Örnek mesaj akışı:

\begin{lstlisting}
TEXTUAL                         BMS_MAIN
   |                               |
   |--- GET_CELL_VOLTAGES -------->|
   |                               |
   |<-- CELL 1 : 3.651 V ----------|
   |                               |
   |       100 ms                  |
   |                               |
   |<-- CELL 2 : 3.648 V ----------|
   |                               |
   |       100 ms                  |
   |                               |
   |<-- CELL 3 : 3.652 V ----------|
   |                               |
   |       ...                     |
   |                               |
   |<-- END_OF_LIST ---------------|
\end{lstlisting}

Liste sonu mesajı:

\begin{lstlisting}
Cell Index   = 0x00
Cell Voltage = 0x0000
\end{lstlisting}

olarak gönderilir.

% =========================================================
\section{Sıcaklık Sensörlerinin Okunması}

\subsection{GET\_TEMPERATURES}

BMS'e bağlı tüm sıcaklık sensörlerinin değerlerini almak için:

\begin{lstlisting}
Message Class = DEOS_CLASS_REQUEST
Service       = DEOS_SERVICE_BMS
Command       = DEOS_CMD_BMS_GET_TEMPERATURES
\end{lstlisting}

kullanılır.

Request command-specific payload içermez.

\subsection{TEMPERATURE}

Her sıcaklık sensörü ayrı bir response frame'i içerisinde gönderilir:

\begin{lstlisting}
Message Class = DEOS_CLASS_RESPONSE
Service       = DEOS_SERVICE_BMS
Command       = DEOS_CMD_BMS_TEMPERATURE
\end{lstlisting}

Command payload:

\begin{center}
\begin{tabular}{c c c l}
\toprule
\textbf{Offset} &
\textbf{Boyut} &
\textbf{Tip} &
\textbf{Alan} \\
\midrule

0 & 1 byte & uint8\_t & Sensor Index \\
1 & 2 byte & int16\_t & Temperature \\

\bottomrule
\end{tabular}
\end{center}

Temperature ölçeği:

\begin{lstlisting}
Physical Temperature = Raw Value * 0.1 degC
\end{lstlisting}

Örnek:

\begin{lstlisting}
Sensor Index = 0x02
Temperature  = 274

-> Sensor 2 = 27.4 degC
\end{lstlisting}

Beş sıcaklık sensörü bulunan bir sistemde örnek akış:

\begin{lstlisting}
GROUND_CONTROL                     BMS_MAIN
      |                                |
      |--- GET_TEMPERATURES ---------->|
      |                                |
      |<-- SENSOR 1 -------------------|
      |          100 ms                |
      |<-- SENSOR 2 -------------------|
      |          100 ms                |
      |<-- SENSOR 3 -------------------|
      |          100 ms                |
      |<-- SENSOR 4 -------------------|
      |          100 ms                |
      |<-- SENSOR 5 -------------------|
      |          100 ms                |
      |<-- END_OF_LIST ----------------|
\end{lstlisting}

Liste sonu:

\begin{lstlisting}
Sensor Index = 0x00
Temperature  = 0x0000
\end{lstlisting}

ile belirtilir.

% =========================================================
\section{Hücre Balancing Durumunun Okunması}

Hücre balancing bilgisi de hücre sayısına bağlı olduğundan liste
yöntemi kullanılır.

\subsection{GET\_BALANCING}

\begin{lstlisting}
Message Class = DEOS_CLASS_REQUEST
Service       = DEOS_SERVICE_BMS
Command       = DEOS_CMD_BMS_GET_BALANCING
\end{lstlisting}

Request payload bulunmaz.

\subsection{BALANCING}

Her hücrenin balancing durumu ayrı bir response ile gönderilir.

\begin{lstlisting}
Message Class = DEOS_CLASS_RESPONSE
Service       = DEOS_SERVICE_BMS
Command       = DEOS_CMD_BMS_BALANCING
\end{lstlisting}

Payload:

\begin{center}
\begin{tabular}{c c c l}
\toprule
\textbf{Offset} &
\textbf{Boyut} &
\textbf{Tip} &
\textbf{Alan} \\
\midrule

0 & 1 byte & uint8\_t & Cell Index \\
1 & 1 byte & uint8\_t & Balancing State \\

\bottomrule
\end{tabular}
\end{center}

Balancing State:

\begin{center}
\begin{tabular}{c l}
\toprule
\textbf{Değer} & \textbf{Anlam} \\
\midrule
0x00 & Balancing inactive \\
0x01 & Balancing active \\
\bottomrule
\end{tabular}
\end{center}

Liste sonunda:

\begin{lstlisting}
Cell Index      = 0x00
Balancing State = 0x00
\end{lstlisting}

gönderilir.

% =========================================================
\section{SYSTEM Service}

BMS alt sisteminde SYSTEM servisi:

\begin{lstlisting}
DEOS_SERVICE_SYSTEM = 0x00
\end{lstlisting}

olarak kullanılır.

BMS\_MAIN, BMS\_AUX ve BMS\_GATEWAY SYSTEM servisini destekler.

Bu alt sistemde kullanılan SYSTEM komutları:

\begin{center}
\begin{tabular}{c l}
\toprule
\textbf{ID} & \textbf{Command} \\
\midrule

0x05 & DEOS\_CMD\_SYSTEM\_RESET\_NODE \\
0x06 & DEOS\_CMD\_SYSTEM\_GET\_NODE\_INFO \\
0x07 & DEOS\_CMD\_SYSTEM\_HEARTBEAT \\
0x08 & DEOS\_CMD\_SYSTEM\_PING \\
0x09 & DEOS\_CMD\_SYSTEM\_PONG \\

\bottomrule
\end{tabular}
\end{center}

Bu komutların genel davranışı ana ICD içerisinde tanımlanmıştır.

GET\_NODE\_INFO response command payload'u:

\begin{center}
\begin{tabular}{c c l}
\toprule
\textbf{Offset} & \textbf{Boyut} & \textbf{Alan} \\
\midrule

0 & 1 byte & Firmware Version \\
1 & 1 byte & Supported Protocol Version \\
2 & 1 byte & Device Type \\

\bottomrule
\end{tabular}
\end{center}

Device Type:

\begin{center}
\begin{tabular}{c l}
\toprule
\textbf{Değer} & \textbf{Anlam} \\
\midrule
0x00 & Invalid / Unknown \\
0x01 & Physical Node \\
0x02 & Virtual Node \\
\bottomrule
\end{tabular}
\end{center}

% =========================================================
\section{DIAGNOSTIC Service}

Diagnostic Service:

\begin{lstlisting}
DEOS_SERVICE_DIAGNOSTIC = 0x06
\end{lstlisting}

olarak kullanılır.

BMS\_MAIN, BMS\_AUX ve BMS\_GATEWAY DIAGNOSTIC servisini destekler.

Kullanılan komutlar:

\begin{center}
\begin{tabular}{c l}
\toprule
\textbf{ID} & \textbf{Command} \\
\midrule

0x01 & DEOS\_CMD\_DIAG\_GET\_HEALTH \\
0x02 & DEOS\_CMD\_DIAG\_GET\_FAULTS \\
0x03 & DEOS\_CMD\_DIAG\_CLEAR\_FAULTS \\

\bottomrule
\end{tabular}
\end{center}

Fault kayıt yapısı, Fault ID namespace'i, Severity, State,
GET\_FAULTS aktarım mekanizması, 100 ms fault aktarım periyodu,
End Of Fault List ve CLEAR\_FAULTS davranışı ayrı
\textit{DEOS Fault Management} dokümanında tanımlanmıştır.

Bu nedenle fault protokolü bu dokümanda tekrar tanımlanmaz.

BMS'e ait fault bilgisi isteniyorsa Destination ilgili BMS node'u
olmalıdır.

Örneğin:

\begin{lstlisting}
Source        = DEOS_NODE_TEXTUAL
Destination   = DEOS_NODE_BMS_MAIN
Message Class = DEOS_CLASS_REQUEST
Service       = DEOS_SERVICE_DIAGNOSTIC
Command       = DEOS_CMD_DIAG_GET_FAULTS
\end{lstlisting}

Gateway'in kendi fault bilgisi için ise:

\begin{lstlisting}
Destination = DEOS_NODE_BMS_GATEWAY
\end{lstlisting}

kullanılır.

% =========================================================
\section{Örnek Kullanım}

\subsection{Yer Kontrol Tarafından Genel BMS Bilgisinin Alınması}

\begin{lstlisting}
Source        = DEOS_NODE_GROUND_CONTROL
Destination   = DEOS_NODE_BMS_MAIN
Message Class = DEOS_CLASS_REQUEST
Service       = DEOS_SERVICE_BMS
Command       = DEOS_CMD_BMS_GET_STATUS
\end{lstlisting}

BMS:

\begin{lstlisting}
Source        = DEOS_NODE_BMS_MAIN
Destination   = DEOS_NODE_GROUND_CONTROL
Message Class = DEOS_CLASS_RESPONSE
Service       = DEOS_SERVICE_BMS
Command       = DEOS_CMD_BMS_STATUS
\end{lstlisting}

cevabını üretir.

\subsection{Textual Node Tarafından Hücre Voltajlarının Alınması}

\begin{lstlisting}
Source        = DEOS_NODE_TEXTUAL
Destination   = DEOS_NODE_BMS_AUX
Message Class = DEOS_CLASS_REQUEST
Service       = DEOS_SERVICE_BMS
Command       = DEOS_CMD_BMS_GET_CELL_VOLTAGES
\end{lstlisting}

BMS\_AUX tüm hücrelerini birer frame halinde 100 ms aralıklarla
gönderir ve aktarımı END\_OF\_LIST mesajıyla sonlandırır.

% =========================================================
\section{Komut Özeti}

\begin{center}
\begin{tabular}{c c l}
\toprule
\textbf{Service} & \textbf{ID} & \textbf{Command} \\
\midrule

BMS & 0x01 & GET\_STATUS \\
BMS & 0x02 & STATUS \\
BMS & 0x03 & GET\_CELL\_VOLTAGES \\
BMS & 0x04 & CELL\_VOLTAGE \\
BMS & 0x05 & GET\_TEMPERATURES \\
BMS & 0x06 & TEMPERATURE \\
BMS & 0x07 & GET\_BALANCING \\
BMS & 0x08 & BALANCING \\

\midrule

SYSTEM & 0x05 & RESET\_NODE \\
SYSTEM & 0x06 & GET\_NODE\_INFO \\
SYSTEM & 0x07 & HEARTBEAT \\
SYSTEM & 0x08 & PING \\
SYSTEM & 0x09 & PONG \\

\midrule

DIAGNOSTIC & 0x01 & GET\_HEALTH \\
DIAGNOSTIC & 0x02 & GET\_FAULTS \\
DIAGNOSTIC & 0x03 & CLEAR\_FAULTS \\

\bottomrule
\end{tabular}
\end{center}

% =========================================================
\section{Uygulama Kuralları}

\begin{enumerate}[nosep]

    \item BMS\_MAIN ve BMS\_AUX bağımsız DEOS node'larıdır.

    \item BMS Gateway yalnızca SYSTEM ve DIAGNOSTIC servislerini
    kendi adına destekler.

    \item BMS Service yalnızca BMS\_MAIN ve BMS\_AUX için geçerlidir.

    \item Bir BMS'ten veri isteyen node doğrudan ilgili mantıksal
    BMS Node ID'sini Destination olarak kullanmalıdır.

    \item GET\_STATUS, sık kullanılan genel BMS bilgilerinin tek
    mesajda alınması için kullanılır.

    \item Hücre voltajları, sıcaklık sensörleri ve balancing durumları
    ayrı liste komutları ile alınır.

    \item Liste tipi transferlerde her response frame'i yalnızca
    tek bir eleman taşır.

    \item Ardışık liste elemanları nominal 100 ms aralıkla gönderilir.

    \item Index = 0x00 liste sonunu ifade eder ve gerçek bir eleman
    numarası olarak kullanılmaz.

    \item Gateway üzerinden taşınan mesajlarda mantıksal Source ve
    Destination Node kimlikleri korunur.

    \item Fault işlemlerinde DEOS Fault Management dokümanı
    uygulanmalıdır.

\end{enumerate}

\end{document}